\documentclass[12pt,a4paper]{article}
\usepackage[margin=20mm]{geometry}
\usepackage[T1]{fontenc}
\usepackage[italian]{babel}
\usepackage{lmodern,amsmath,amssymb,tikz,xcolor}
\definecolor{ink}{HTML}{182536}
\definecolor{accent}{HTML}{085872}
\definecolor{rulecolor}{HTML}{7C8B95}
\definecolor{light}{HTML}{EEF4F6}
\setlength{\parindent}{0pt}
\pagestyle{empty}
\color{ink}
\newcommand{\heading}[1]{%
  {\large\sffamily\bfseries\color{accent} MATEMATICA}\par\vspace{2mm}
  {\Huge\sffamily\bfseries Sistema n.~#1}\par\vspace{3mm}
  {\color{accent}\rule{\linewidth}{1pt}}\par}
\newcommand{\step}[1]{\par\vspace{3mm}{\large\sffamily\bfseries\color{accent}#1}\par\vspace{1mm}}
\newcommand{\dotat}[3]{\draw[ink,fill=#3,line width=.8pt] (#1,#2) circle[radius=2.5pt];}
\newcommand{\circledsign}[3]{\node[draw=accent,fill=light,circle,line width=1pt,minimum size=5.5mm,inner sep=0pt] at (#1,#2) {$#3$};}
\newcommand{\plain}[3]{\node at (#1,#2) {$#3$};}
\newcommand{\signframe}[2]{%
  \draw[fill=light,draw=none] (0,-2.75) rectangle (14,-1.95);
  \foreach \y in {-.4,-1.17,-1.95,-2.75}
    \draw[rulecolor,thin] (0,\y)--(14,\y);
  \foreach \x in {3,6,10}
    \draw[rulecolor,thin,densely dotted] (\x,-.4)--(\x,-2.75);
  \node[anchor=west] at (.2,0) {$x$};
  \plain{6}{0}{#1}
  \plain{10}{0}{#2}
  \node[anchor=west] at (.25,-.79) {$A$};
  \node[anchor=west] at (.25,-1.56) {$B$};
  \node[anchor=west,align=left] at (.25,-2.35) {$A/B$};

}
\newcommand{\systems}[3]{%
  \draw[->,rulecolor] (3,0)--(14,0);
  \foreach \x/\labelx in {5/#1,8/#2,11/#3} {
    \draw[rulecolor,thin,densely dotted] (\x,-.07)--(\x,-2.98);
    \node[above] at (\x,0) {$\labelx$};
  }
  \node[anchor=west] at (.15,-.76) {$S_1$};
  \node[anchor=west] at (.15,-1.60) {$S_2$};
  \node[anchor=west] at (.15,-2.56) {$S_1\cap S_2$};
}
\newcommand{\on}[3]{\draw[ink,line width=1.35pt] (#1,#3)--(#2,#3);}
\newcommand{\off}[3]{\draw[rulecolor,line width=.8pt,dash pattern=on 3pt off 2.5pt] (#1,#3)--(#2,#3);}
\newcommand{\legend}{\par\smallskip{\small Linea continua per le soluzioni, tratteggiata per il complementare.}\par}
\newcommand{\sol}[2]{%
  \par\vspace{3mm}\noindent\colorbox{light}{\parbox{\dimexpr\linewidth-2\fboxsep}{%
  \centering\vspace{1mm}{\sffamily\bfseries\color{accent} Soluzione del sistema}\hspace{8mm}%
  $\displaystyle S=#1\quad\Longleftrightarrow\quad #2$\par\vspace{1mm}}}}
\begin{document}
\heading{170}
\[
\begin{cases}
\dfrac{x-4}{x+1}\ge 0\\[3pt]
-4(2-x)\le 0
\end{cases}
\]
\step{Prima disequazione}
Poniamo $A=x-4$ e $B=x+1$. La condizione di esistenza è $x\ne -1$.
Per lo studio dei segni poniamo $A\ge 0$ e $B>0$.
\[
\begin{aligned}
A\ge 0 &\quad\Longleftrightarrow\quad x-4\ge 0\quad\Longleftrightarrow\quad x\ge 4,\\
B>0 &\quad\Longleftrightarrow\quad x+1>0\quad\Longleftrightarrow\quad x>-1.
\end{aligned}
\]
\begin{center}
\begin{tikzpicture}[x=1.03cm,y=.83cm]
  \signframe{-1}{4}
  \plain{4.5}{-.79}{-}\plain{8}{-.79}{-}\plain{12}{-.79}{+}
  \plain{4.5}{-1.56}{-}\plain{8}{-1.56}{+}\plain{12}{-1.56}{+}
  \dotat{10}{-.79}{ink}\dotat{6}{-1.56}{white}
  \circledsign{4.5}{-2.35}{+}\plain{8}{-2.35}{-}\circledsign{12}{-2.35}{+}
  \circledsign{10}{-2.35}{0}
  \node[font=\scriptsize\sffamily,fill=light,inner sep=1pt] at (6,-2.35) {escluso};
\end{tikzpicture}
\end{center}
Cerchiamo i $+$ e lo $0$, perché il quoziente deve essere $\ge 0$.
Il valore $-1$ è escluso perché annulla il denominatore.
\[
S_1=(-\infty,-1)\cup[4,+\infty).
\]
\step{Seconda disequazione}
Dividendo per $-4$, si inverte il verso della disequazione.
\[
-4(2-x)\le 0\quad\Longleftrightarrow\quad 2-x\ge 0
\quad\Longleftrightarrow\quad x\le 2,
\qquad S_2=(-\infty,2].
\]
\step{Intersezione delle soluzioni}
\legend
\begin{center}
\begin{tikzpicture}[x=1.03cm,y=.83cm]
  \systems{-1}{2}{4}
  \on{3}{5}{-.76}\off{5}{11}{-.76}\on{11}{14}{-.76}
  \dotat{5}{-.76}{white}\dotat{11}{-.76}{ink}
  \on{3}{8}{-1.60}\off{8}{14}{-1.60}\dotat{8}{-1.60}{ink}
  \on{3}{5}{-2.56}\off{5}{14}{-2.56}\dotat{5}{-2.56}{white}
\end{tikzpicture}
\end{center}
\sol{(-\infty,-1)}{x<-1}
\newpage
\heading{171}
\[
\begin{cases}
\dfrac{x-1}{4-x}\le 0\\[6pt]
\dfrac{x-3}{4}<\dfrac{x}{6}
\end{cases}
\]
\step{Prima disequazione}
Poniamo $A=x-1$ e $B=4-x$. La condizione di esistenza è $x\ne 4$.
Per il denominatore usiamo $B>0$ e invertiamo il verso quando isoliamo $x$.
\[
\begin{aligned}
A\ge 0 &\quad\Longleftrightarrow\quad x-1\ge 0\quad\Longleftrightarrow\quad x\ge 1,\\
B>0 &\quad\Longleftrightarrow\quad 4-x>0\quad\Longleftrightarrow\quad x<4.
\end{aligned}
\]
\begin{center}
\begin{tikzpicture}[x=1.03cm,y=.83cm]
  \signframe{1}{4}
  \plain{4.5}{-.79}{-}\plain{8}{-.79}{+}\plain{12}{-.79}{+}
  \plain{4.5}{-1.56}{+}\plain{8}{-1.56}{+}\plain{12}{-1.56}{-}
  \dotat{6}{-.79}{ink}\dotat{10}{-1.56}{white}
  \circledsign{4.5}{-2.35}{-}\plain{8}{-2.35}{+}\circledsign{12}{-2.35}{-}
  \circledsign{6}{-2.35}{0}
  \node[font=\scriptsize\sffamily,fill=light,inner sep=1pt] at (10,-2.35) {escluso};
\end{tikzpicture}
\end{center}
Cerchiamo i $-$ e lo $0$, perché il quoziente deve essere $\le 0$.
Il valore $4$ è escluso perché annulla il denominatore.
\[
S_1=(-\infty,1]\cup(4,+\infty).
\]
\step{Seconda disequazione}
Moltiplichiamo entrambi i membri per $12$, il minimo comune multiplo dei denominatori.
\[
\dfrac{x-3}{4}<\dfrac{x}{6}
\quad\Longleftrightarrow\quad 3(x-3)<2x
\quad\Longleftrightarrow\quad x<9,
\qquad S_2=(-\infty,9).
\]
\step{Intersezione delle soluzioni}
\legend
\begin{center}
\begin{tikzpicture}[x=1.03cm,y=.83cm]
  \systems{1}{4}{9}
  \on{3}{5}{-.76}\off{5}{8}{-.76}\on{8}{14}{-.76}
  \dotat{5}{-.76}{ink}\dotat{8}{-.76}{white}
  \on{3}{11}{-1.60}\off{11}{14}{-1.60}\dotat{11}{-1.60}{white}
  \on{3}{5}{-2.56}\off{5}{8}{-2.56}\on{8}{11}{-2.56}\off{11}{14}{-2.56}
  \dotat{5}{-2.56}{ink}\dotat{8}{-2.56}{white}\dotat{11}{-2.56}{white}
\end{tikzpicture}
\end{center}
\sol{(-\infty,1]\cup(4,9)}{x\le 1\ \lor\ 4<x<9}
\end{document}
